\documentclass[12pt,a4paper]{article}

\usepackage{amsmath}
\usepackage{amssymb}
\usepackage{xeCJK}
\usepackage{geometry}
\usepackage{graphicx}
\geometry{margin=2.5cm}

\setCJKmainfont{Noto Sans CJK TC}

\setlength{\parindent}{0pt}
\setlength{\parskip}{6pt}

\begin{document}

\section{Derivation: $V_{out}/V_{in}$}

\begin{enumerate}
    \item 令輸入為 $I_{in}$,輸出為 $V_{out} = V_{ij}$。

    其中
    \begin{equation}
        V_{ij} = S \cdot B_{ij}
    \end{equation}

    則
    \begin{equation}
        \frac{V_{ij}}{I_{in}} = \frac{S}{I_{in}}
        \begin{bmatrix}
            B_{11} & \cdots & \cdots & \cdots & \cdots & B_{16} \\
            B_{21} & \cdots & \cdots & \cdots & \cdots & B_{26} \\
            B_{31} & \cdots & \cdots & \cdots & \cdots & \vdots \\
            B_{41} & \cdots & \cdots & \cdots & \cdots & \vdots \\
            B_{51} & \cdots & \cdots & \cdots & \cdots & \vdots \\
            B_{61} & \cdots & \cdots & \cdots & \cdots & B_{66}
        \end{bmatrix}
        = \frac{S}{I_{in}} \overline{B}_S
    \end{equation}

    其中 $B_{ij}$ 為激發第 $j$ 個磁極、從第 $i$ 個磁極量測到之磁通密度,$S$ 為霍爾感測器敏感度。

    \item 考慮放大器 $\Rightarrow I_{in} = V_{in} \cdot k_A$

    則
    \begin{equation}
        \frac{V_{out}}{V_{in}} = \underbrace{\frac{k_A \cdot S}{I_{in}}}_{\alpha}\,\overline{B}_S
    \end{equation}
\end{enumerate}

\section{Sensor placement}

每顆磁極都依照 Long Fei 標準:沿磁極表面 4.572 mm + 垂直磁極表面 0.41 mm,Hall sensor 直徑 0.3 mm,disc 平均(721 點 annular sampling)。

\textbf{Lower poles (P1/P3/P6)} — milled flat (halfcut 切割面),$\hat{n}_+$ 純垂直 milled flat 向外。

\textbf{Upper poles (P2/P4/P5)} — 沿 cone top slant 4.572 mm + 垂直 slant 0.41 mm,$\hat{n}_+ = (-\sin\beta, \cos\beta, 0)$ in pole-local,$\beta = \arctan(3/15) \approx 11.31^\circ$。

\section{FEM 結果 (Long2016 verbatim baseline)}

\textbf{模擬條件}:Long2016 verbatim halfcut hexapole,$\mu_r = 280$,$I_{in} = 1.0$ A,每顆 coil 70 匝 CCW (looking from $+z$),mesh node $\approx 494{,}873$ 一致。Sign convention: 對每根磁極 $j$,raw FEM 矩陣的第 $j$ 行乘以 $\mathrm{sign}(\bar{B}_S(j,j))$ 使對角為正(+I 對應 yoke $\to$ tip $\to$ WP)。

\subsection{$\overline{B}_S$ matrix (Tesla)}

\[
\bar{B}_S =
\begin{bmatrix}
 +0.010611 &  -0.000414 &  -0.000163 &  -0.004500 &  -0.004360 &  -0.000186 \\
 +0.000060 &  +0.009543 &  -0.001660 &  -0.002845 &  -0.002844 &  -0.001662 \\
 -0.000238 &  -0.004803 &  +0.011469 &  -0.000450 &  -0.004619 &  -0.000268 \\
 -0.001596 &  -0.002807 &  +0.000075 &  +0.009429 &  -0.002838 &  -0.001586 \\
 -0.001642 &  -0.002828 &  -0.001588 &  -0.002846 &  +0.009537 &  +0.000052 \\
 -0.000167 &  -0.004186 &  -0.000137 &  -0.004214 &  -0.000403 &  +0.010156
\end{bmatrix}\ \mathrm{T}
\]


\subsection{$V_{out}/V_{in}$ matrix (dimensionless)}

$V_{out}/V_{in} = k_A \cdot S_{\mathrm{Hall}} / I_{in} \cdot \overline{B}_S$,with $k_A = 0.3614$\,A/V (power amplifier gain),$S_{\mathrm{Hall}} = 130$\,V/T (EQ-730L)。

\[
\frac{V_{out}}{V_{in}} = \begin{bmatrix}
  +0.4985 &   -0.0195 &   -0.0076 &   -0.2114 &   -0.2048 &   -0.0087 \\
  -0.0028 &   +0.4484 &   -0.0780 &   -0.1337 &   -0.1336 &   -0.0781 \\
  -0.0112 &   -0.2256 &   +0.5389 &   -0.0211 &   -0.2170 &   -0.0126 \\
  -0.0750 &   -0.1319 &   -0.0035 &   +0.4430 &   -0.1333 &   -0.0745 \\
  -0.0771 &   -0.1329 &   -0.0746 &   -0.1337 &   +0.4481 &   -0.0024 \\
  -0.0078 &   -0.1967 &   -0.0064 &   -0.1980 &   -0.0189 &   +0.4771
\end{bmatrix}
\]


\subsection{Lower vs Upper diagonal}

\begin{tabular}{lll}
\hline
組別 & 磁極 & $\bar{B}_S$ 對角 (mT) \\
\hline
Lower (halfcut D-shape) & P1, P3, P6 & 10.6, 11.5, 10.2 (avg 10.74) \\
Upper (full cone)       & P2, P4, P5 & 9.5, 9.4, 9.5 (avg 9.49) \\
\hline
\multicolumn{2}{l}{比例 Lower / Upper} & 1.13 \\
\hline
\end{tabular}

Lower halfcut 把 cone 削掉一半,殘餘截面 flux 集中,fringe 在 milled flat 上方反而更強;upper full cone flux 在整個截面分散,fringe 較弱。

\end{document}
